Datasheets for datasets are intended to address the needs of two key
stakeholder groups: dataset creators and dataset consumers. For
dataset creators, the primary objective is to encourage careful
reflection on the process of creating, distributing, and maintaining a
dataset, including any underlying assumptions, potential risks or
harms, and implications of use. For dataset consumers, the primary
objective is to ensure they have the information they need to make
informed decisions about using a \edit{dataset.} Transparency on the
part \edit{of} dataset creators is necessary for dataset consumers to
be sufficiently well informed that they can select appropriate
datasets for their
\edit{chosen} tasks and avoid unintentional misuse.\footnote{\edit{We note that in
some cases, the people creating a datasheet for a dataset may not be
the dataset creators, as was the case with the example datasheets that
we created as part of our development process.}}

Beyond these two key stakeholder groups, datasheets \edit{for datasets} may
\edit{be} valuable to policy makers, consumer advocates, \edit{investigative
journalists}, individuals whose data is included in datasets, and
\edit{individuals} who may be impacted by models trained or evaluated \edit{using} datasets. They also serve a secondary objective of facilitating
greater reproducibility of machine learning results: researchers and
practitioners \edit{without access to a dataset may be able to use the
information in its datasheet to create alternative datasets with
similar characteristics.}

Although we provide a set of questions designed to elicit the
information that a datasheet for a dataset might contain, \edit{these
questions} are not intended to be prescriptive. Indeed, we expect that
datasheets will \edit{necessarily} vary depending on factors such as the
domain or existing organizational infrastructure and workflows. For
example,
\edit{some the questions are appropriate for academic researchers publicly
releasing datasets for the purpose of enabling future research, but
less relevant for product teams creating internal datasets for
training proprietary models.  As another example,}~\citet{statements}
outline a proposal similar to datasheets for
datasets \edit{specifically intended for language-based
datasets. Their} questions may \edit{be} naturally integrated into a
datasheet for a language-based dataset as appropriate.

We emphasize that the process of creating a datasheet is not intended
to be automated. Although automated documentation processes are
convenient, they run counter to our objective of encouraging dataset
creators to carefully reflect on the process of creating,
distributing, and maintaining a dataset.
